% !PRINTED 473
\weberpage{473}
%<<CONT
\weq{26}{\frac{c+\partial\omega}{a} = \frac{2\gamma+\delta\omega}{2\alpha+\beta\omega}, \quad \delta \equiv 0 \pmod{2}}

und außerdem

\weq{27}{\gamma^2 - 1 + \gamma(2\alpha-\delta) \equiv 0 \pmod{16};}

aus (26) und (24) folgt:
\[
\begin{aligned}
2\alpha &= B+y, & 2n\gamma &\equiv -C\\
\beta &= A, & 2n\delta &\equiv -(B-y)
\end{aligned}
\quad (\mathrm{mod}\ 16),
\]
und daraus schließt man wie oben:

9. \gesperrt{Ist $m \equiv 2 \pmod{4}$, so ist $\sqrt{2}\,f_1(\omega)^6$ Klasseninvariante.}

Wiederum weisen sämtliche Beispiele darauf hin, daß, wenn
$m \equiv 4 \pmod{8}$ ist, $\sqrt{2}\,f_1(\omega)^{12}$ \gesperrt{Klasseninvariante} ist. Aber
auch hierfür fehlt noch der Beweis. Wenn $m$ durch 8 teilbar
ist, tritt eine Reduktion nicht ein.

In den Fällen 4. bis 8. ergibt sich nach 3. eine weitere
Reduktion auf die dritte Wurzel, wenn $m$ nicht durch 3 teilbar
ist. Demnach erhalten wir folgende Fälle:
\[
m \equiv \pm 1 \pmod{3},
\]
\[
\begin{array}{lll}
 & & \text{Klasseninvariante}\\
1) & m \equiv 3 \pmod{4} & f(\omega)^2,\\
2) & m \equiv 5 \pmod{8} & f(\omega)^4,\\
3) & m \equiv 1 \ \text{\textquotedbl} & \sqrt{2}\,f(\omega)^2,\\
4) & m \equiv 7 \ \text{\textquotedbl} & \sqrt{2}\,f(\omega),\\
5) & m \equiv 3 \ \text{\textquotedbl} & f(\omega),\\
6) & m \equiv 2 \ \text{\textquotedbl} & \sqrt{2}\,f_1(\omega)^2,\\
7) & m \equiv 4 \ \text{\textquotedbl} & \sqrt{2}\,f_1(\omega)^4,
\end{array}
\]
wovon freilich die Fälle 5) und 7) nur durch Induktion geschlossen
 sind.

Die Klasseninvarianten $f(\omega)$ eignen sich ganz besonders zur
numerischen Berechnung, weil sie unter allen die einfachsten
Resultate liefern. Wir geben daher zunächst eine größere Reihe
von Beispielen, die sich auf Grund unserer bisherigen Entwickelungen leicht ableiten lassen, und die zugleich die Methoden
kennen lehren, deren man sich auch zu weiter fortgesetzten Rechnungen dieser Art bedienen kann.

Es wird bei diesen Berechnungen häufig die Aufgabe gestellt,
reduzible, ganze rationale Funktionen in Faktoren zu zerlegen.
Eine solche Zerlegung ist, wenn sie gefunden ist, natürlich sofort
%CONT>>
